% ============================================================================
% TECHNICAL APPENDIX 8A: HITL System Architecture Patterns
% Formal patterns, pseudocode, and implementation considerations
% ============================================================================
% Source: /chapters/ch8/Ch8A_final.md
% ============================================================================

\chapter*{Technical Appendix 8A: HITL System Architecture Patterns}
\addcontentsline{toc}{chapter}{Technical Appendix 8A: HITL System Architecture Patterns}

\begin{chapterquote}{Formal foundations for Chapter 8}
Chapter~8 introduced the blueprint for smart helper systems: routing, queuing, monitoring, and feedback. This appendix provides formal definitions, complete reference implementations, and production considerations for each component.
\end{chapterquote}

% ============================================================================
\section{The Pipeline vs.\ Loop Distinction}
% ============================================================================

A \term{HITL Pipeline} is a directed acyclic graph $G = (V, E)$ where: $V$ includes exactly one human review node $H$; $H$ receives input from a model node $M$; $H$ produces output consumed by a downstream process; and no path exists from any downstream node back to $M$.

A \term{HITL Loop} is a directed graph $G = (V, E)$ where: $V$ includes at least one human review node $H$; and a path exists from $H$ back to some component that influences future model behaviour.

The loop property is necessary but not sufficient for a learning HITL system. The path from $H$ back to the model must also satisfy:
\begin{enumerate}
    \item Human labels are available for integration into training.
    \item Labels are quality-filtered before integration.
    \item The training process is actually triggered periodically.
\end{enumerate}

\begin{cautionbox}[The Silent Queue]
A common failure mode is a nominally ``looping'' HITL system whose feedback queue is never processed. The system collects human labels, stores them, and retrains on schedule---except the schedule was never activated in production. Verify that the retraining job is running, not just that the queue is filling.
\end{cautionbox}

% ============================================================================
\section{The Three-Zone Router}
% ============================================================================

The core routing component partitions incoming cases into three zones: auto-positive, auto-negative, and human review, with dynamic threshold adjustment based on reviewer capacity.

\begin{lstlisting}[style=python, caption={Three-zone router with capacity adaptation}]
from dataclasses import dataclass
from enum import Enum
from typing import Optional, Callable

class Disposition(Enum):
    AUTO_POSITIVE = "auto_positive"
    AUTO_NEGATIVE = "auto_negative"
    HUMAN_REVIEW  = "human_review"

@dataclass
class RoutingConfig:
    tau_low:           float     # Below this: auto-negative
    tau_high:          float     # Above this: auto-positive
    capacity_fn:       Callable  # Returns current reviewer capacity [0,1]
    overflow_narrowing: float = 0.05
    underflow_widening: float = 0.03

class ThreeZoneRouter:
    def __init__(self, config: RoutingConfig):
        self.config    = config
        self.tau_low   = config.tau_low
        self.tau_high  = config.tau_high

    def route(self, score: float,
              stakes: Optional[float] = None) -> Disposition:
        eff_high = self.tau_high
        eff_low  = self.tau_low
        if stakes is not None and stakes > 0.8:
            narrowing = 0.15
            eff_high = self.tau_high - narrowing
            eff_low  = self.tau_low  + narrowing
        if score >= eff_high:
            return Disposition.AUTO_POSITIVE
        elif score <= eff_low:
            return Disposition.AUTO_NEGATIVE
        else:
            return Disposition.HUMAN_REVIEW

    def adjust_for_capacity(self):
        capacity = self.config.capacity_fn()
        if capacity < 0.3:   # Overloaded: narrow the review band
            self.tau_low  = min(self.tau_low  + self.config.overflow_narrowing, 0.45)
            self.tau_high = max(self.tau_high - self.config.overflow_narrowing, 0.55)
        elif capacity > 0.8: # Underutilised: widen the band
            base_low  = self.config.tau_low
            base_high = self.config.tau_high
            self.tau_low  = max(self.tau_low  - self.config.underflow_widening,
                                base_low  - 0.15)
            self.tau_high = min(self.tau_high + self.config.underflow_widening,
                                base_high + 0.15)
\end{lstlisting}

\begin{examplebox}[Stakes Override]
The \texttt{stakes} parameter implements a high-stakes override: when a case is flagged as high-stakes (e.g., a medical decision, a financial transaction above a threshold), the review band narrows even if the model is moderately confident. This ensures high-consequence cases always pass through human review regardless of reviewer capacity.
\end{examplebox}

% ============================================================================
\section{Active Learning Queue Implementation}
% ============================================================================

Cases routed for human review should be prioritised not by arrival order but by \emph{learning value}: the expected improvement to the model from labelling this particular case. The implementation below combines uncertainty (proximity to the decision boundary) with diversity (dissimilarity from already-queued cases).

\begin{lstlisting}[style=python, caption={Active learning priority queue}]
from dataclasses import dataclass, field
from typing import List
import numpy as np
from heapq import heappush, heappop, heapify, nsmallest

@dataclass(order=True)
class LearningQueueItem:
    learning_value: float
    score:    float      = field(compare=False)
    case_id:  str        = field(compare=False)
    features: np.ndarray = field(compare=False)

    def __post_init__(self):
        object.__setattr__(self, 'learning_value',
                           -abs(self.learning_value))  # max-heap

class ActiveLearningQueue:
    def __init__(self, uncertainty_weight=0.7,
                 diversity_weight=0.3, max_size=1000):
        self.heap = []
        self.uw   = uncertainty_weight
        self.dw   = diversity_weight
        self.max_size = max_size
        self._labeled = []

    def _uncertainty(self, score: float) -> float:
        return 1 - abs(2 * score - 1)  # peaks at 0.5

    def _diversity(self, features: np.ndarray) -> float:
        if not self._labeled:
            return 1.0
        labeled   = np.array(self._labeled)
        min_dist  = np.linalg.norm(labeled - features, axis=1).min()
        return min(min_dist / (min_dist + 1.0), 1.0)

    def add(self, case_id: str, score: float, features: np.ndarray):
        lv = self.uw * self._uncertainty(score) + \
             self.dw * self._diversity(features)
        heappush(self.heap, LearningQueueItem(lv, score,
                                              case_id, features))
        if len(self.heap) > self.max_size:
            self.heap = nsmallest(self.max_size, self.heap)
            heapify(self.heap)

    def pop_top_k(self, k: int) -> List[LearningQueueItem]:
        results = []
        for _ in range(min(k, len(self.heap))):
            item = heappop(self.heap)
            results.append(item)
            self._labeled.append(item.features)
        return results
\end{lstlisting}

% ============================================================================
\section{Production Monitoring}
% ============================================================================

A production HITL system requires continuous monitoring of three signals: \emph{confidence distribution drift} (is the model seeing different inputs than it was trained on?), \emph{human override rate} (are reviewers frequently disagreeing with automated decisions?), and \emph{review rate} (is the fraction of cases reaching human review unexpectedly high?).

\begin{lstlisting}[style=python, caption={HITL system health monitor}]
from collections import deque
from datetime import datetime
from typing import Dict, Optional
import numpy as np
from scipy import stats

class HITLMonitor:
    def __init__(self, reference_confidence_dist: np.ndarray,
                 window_size: int = 1000,
                 review_band: tuple = (0.3, 0.7),
                 alert_thresholds: Optional[Dict] = None):
        self.reference = reference_confidence_dist
        self.conf_window     = deque(maxlen=window_size)
        self.override_window = deque(maxlen=window_size)
        self.review_low, self.review_high = review_band
        self.thresholds = alert_thresholds or {
            'drift_pvalue':        0.01,
            'override_rate_high':  0.30,
            'review_rate_high':    0.60,
        }

    def record(self, score: float, disposition: str,
               human_override: Optional[bool] = None):
        self.conf_window.append(score)
        if human_override is not None:
            self.override_window.append(int(human_override))

    def check_health(self) -> Dict:
        if len(self.conf_window) < self.conf_window.maxlen // 2:
            return {'status': 'insufficient_data'}
        current = np.array(self.conf_window)
        ks_stat, p_value = stats.ks_2samp(self.reference, current)
        issues = {}
        if p_value < self.thresholds['drift_pvalue']:
            issues['distribution_drift'] = {
                'severity': 'HIGH',
                'ks_stat':  round(ks_stat, 4),
                'p_value':  round(p_value, 6),
            }
        if len(self.override_window) >= 50:
            override_rate = np.mean(list(self.override_window))
            if override_rate > self.thresholds['override_rate_high']:
                issues['high_override_rate'] = {
                    'severity':      'MEDIUM',
                    'override_rate': round(override_rate, 3),
                }
        review_rate = sum(self.review_low <= s <= self.review_high
                          for s in self.conf_window) \
                      / len(self.conf_window)
        if review_rate > self.thresholds['review_rate_high']:
            issues['high_review_rate'] = {
                'severity':    'LOW',
                'review_rate': round(review_rate, 3),
            }
        return {'status': 'issues_detected' if issues else 'healthy',
                'issues': issues}
\end{lstlisting}

\begin{cautionbox}[KS Test Limitations]
The Kolmogorov--Smirnov test used above detects distributional drift in one-dimensional confidence scores. It does not detect drift in the high-dimensional feature space that produces those scores. A model may show stable confidence distributions while operating on inputs that have shifted significantly from its training distribution. Complement confidence monitoring with feature-space monitoring using dimensionality reduction.
\end{cautionbox}

% ============================================================================
\section{Exercises}
% ============================================================================

\begin{Exercise}[Feedback Integration Layer]
Implement a feedback integration layer that: ingests human labels from the review queue; quality-filters them using inter-annotator agreement scores; and triggers retraining when accumulated new labels exceed a threshold. What safeguards prevent bad labels from degrading the model?
\end{Exercise}

\begin{Exercise}[Time-of-Day Threshold Adjustment]
Extend the \texttt{ThreeZoneRouter} to implement time-of-day threshold adjustment: lower the review band during business hours (more capacity) and narrow it during nights and weekends. How would you test whether this adjustment improves overall system performance?
\end{Exercise}

\begin{Exercise}[Ghost Reviewer Detection]
Implement a detector for the ``ghost reviewer'' failure mode: a review queue that has been accumulating items without being processed. What metrics would indicate this? How would you design an alert that distinguishes reviewer absence from legitimate queue growth?
\end{Exercise}

\begin{Exercise}[High-Dimensional Drift Detection]
The monitor checks for distribution drift using a KS test on confidence scores. What are the limitations of this approach? Propose an alternative using dimensionality reduction of input features (e.g., UMAP or t-SNE on penultimate-layer activations). How would you set alert thresholds?
\end{Exercise}

\bigskip
\begin{center}
\rule{0.5\textwidth}{0.4pt}
\end{center}
\bigskip

\noindent\textit{This appendix provides reference implementations for the core components of a production HITL architecture: routing, active learning queues, and system health monitoring, as introduced in Chapter~8.}
